\section{Introduction}
\label{sec:intro}

The paradigm of serving specialized deep learning capabilities has undergone a profound transformation. Rather than training and deploying gargantuan, monolithic models for every individual task, modern machine learning increasingly relies on modular Parameter-Efficient Fine-Tuning (PEFT), most notably Low-Rank Adaptation (LoRA) \cite{hu2021lora}. In dynamic serving environments, multiple task-specific LoRA adapters are served on top of a single pre-trained backbone, and a dynamic router determines how these adapters are blended or ensembled on a query-by-query basis \cite{sable, lora-hub, mofte}.

This dynamic ensembling problem is particularly acute in sequential streaming contexts, where the model encounters a continuous stream of queries belonging to diverse, often shifting, task distributions \cite{chemmerge, packinetics}. When serving such streams, the routing system must manage a fundamental trade-off: \emph{responsiveness} (the ability to instantly transition to a new expert when a task boundary is crossed) versus \emph{stability} (the ability to resist query-level activation noise and maintain coherent representation trajectories across depth and time).

Existing routing methodologies can be split into two main paradigms:
\begin{enumerate}
    \item \textbf{Stateless Routers:} Methods such as SABLE \cite{sable} or classic Mixture-of-Experts (MoE) routing \cite{shazeer2017outrageously, fedus2022switch} project the activations of each query independently to determine routing weights. While highly responsive, they suffer from severe temporal query-to-query routing jitter under sequence-level activation noise. Furthermore, because they lack any active depth-wise recurrence, their ensembling weights remain extremely soft and close to uniform across all network layers, leading to severe representation leakage and parameter co-dominance under overlapping task distributions.
    \item \textbf{Stateful Routers:} Recognizing the jitter paradox, stateful systems attempt to smooth trajectories across layers or time. ChemMerge \cite{chemmerge} models continuous-time biochemical reactions governed by ordinary differential equations (ODEs). While conceptually intriguing, continuous-time solvers are computationally heavy, require ad-hoc numerical clamping hacks to keep weights on the probability simplex, and are highly sensitive to virtual-time hyperparameter configurations. Momentum-Merge \cite{momentummerge} simplifies this via a constant Exponential Moving Average (EMA) of routing weights across depth. More recently, PAC-Kinetics \cite{packinetics} formulates a learned linear stateful router optimized via a PAC-Bayesian generalization bound.
\end{enumerate}

Despite their mathematical tractability, all existing stateful systems (ChemMerge, Momentum-Merge, PAC-Kinetics) are fundamentally constrained by a \textbf{linear state-space assumption}. They model expert interaction as simple decay-and-injection mechanics ($s_t = A s_{t-1} + W e_t$). This linear framework is structurally incapable of modeling the non-linear, multi-stable, and self-regulating competitive dynamics that naturally arise in complex natural systems. As a result, they suffer from a severe \emph{representational lag} (phase delay) during rapid, heterogeneous task switches, or require fragile hand-tuned coefficients to prevent co-dominance and representation interference.

\begin{figure}[t]
    \centering
    \includegraphics[width=\columnwidth]{results/fig1.png}
    \caption{\textbf{Joint Serving Accuracy Comparison on Orthogonal Manifolds.} Our proposed Lotka-Volterra Competitive Serving (LVCS) achieves superior performance across both homogeneous and heterogeneous streams, resolving the representational lag that degrades other stateful methods under rapid task switching.}
    \label{fig:teaser}
\end{figure}

In this paper, we reject the linear state-space paradigm. Drawing inspiration from mathematical ecology and complex systems theory, we propose a novel framework: \textbf{Lotka-Volterra Competitive Serving (LVCS)}. We treat the latent representation space of the deep network as a localized ecosystem, where different task-specific LoRA adapters behave like distinct biological species competing for "resources" (quantified by activation similarity coordinate projections).

Instead of independent linear decay, the virtual population (routing state) of each expert species evolves across network layers governed by a discrete-time \textbf{Lotka-Volterra Ricker competition model} \cite{ricker1954stock, may1976simple}. This formulation guarantees several highly desirable mathematical properties:
\begin{itemize}
    \item \textbf{Guaranteed Positivity:} The exponential form of the Ricker equation mathematically guarantees that expert populations remain strictly positive across all layers, completely bypassing the ad-hoc "clamping hacks" that plague continuous ODE solvers.
    \item \textbf{Non-Linear Competitive Saturation:} The coupled non-linear terms naturally model carrying capacities and inter-species competition, enabling the routing system to learn sharp, multi-stable decision boundaries and suppress high-frequency noise.
    \item \textbf{Adaptive Niche Plasticity:} To eliminate the representational lag during sudden task switches, we scale the inter-species competition coefficients by the local temporal homogeneity of the query stream. Under a sudden transition, the local similarity drops, suspending inter-species competition and allowing the colonizing species (the new task expert) to rapidly establish dominance without historical drag.
\end{itemize}

Our contributions can be summarized as follows:
\begin{enumerate}
    \item We propose \textbf{LVCS}, the first biologically-grounded, non-linear stateful router for parameter-efficient expert serving, bridging discrete population ecology and deep representation learning.
    \item We formulate the \textbf{Lotka-Volterra Ricker Recurrence} for layer-by-layer ensembling, proving that it guarantees strict population positivity and natural self-regulation.
    \item We introduce \textbf{Adaptive Niche Plasticity}, a stream-homogeneity-gated competition mechanism that successfully eliminates phase delay under rapid task transitions.
    \item We design a \textbf{Systems-First Static Coordinate Approximation} that extracts routing coordinates once at an early layer, reducing serving latency by over 51\% compared to a fully dynamic model while maintaining virtually identical ensembling accuracy.
    \item We introduce standard unconstrained non-linear recurrent router baselines (such as a spatially recurrent GRU Router) to deconstruct the trade-offs in dynamic ensembling. We reveal that while unconstrained models are highly expressive, they lack mathematical bounds and exhibit up to \textbf{2$\times$ higher layer-to-layer routing jitter} compared to LVCS.
    \item Through exhaustive evaluation over five random seeds under Orthogonal and Overlapping manifold structures, we show that LVCS achieves outstanding results. On overlapping manifolds, LVCS outperforms the state-of-the-art PAC-Kinetics baseline by up to \textbf{+1.38\%} absolute accuracy (89.08\% vs. 88.00\% on homogeneous, and 90.06\% vs. 88.68\% on heterogeneous streams) under static coordinate serving, and up to \textbf{+1.54\%} under dynamic serving, while dramatically reducing routing jitter compared to unconstrained recurrent baselines.
\end{enumerate}
